\chapter{Principio de Acotación Uniforme}

\begin{definicion}
    Sean $X, Y$ espacios normados y sea $\{T_\alpha\}_{\alpha \in A} \subseteq BL(X, Y)$. Se dice que $\{T_\alpha\}_{\alpha \in A}$ es \textbf{uniformemente acotada} si
    \begin{equation*}
        \exists M > 0 : \|T_\alpha\| \leq M \quad \forall \alpha \in A
    \end{equation*}
\end{definicion}

\begin{definicion}
    Sean $X, Y$ espacios normados y sea $\{T_\alpha\}_{\alpha \in A} \subseteq BL(X, Y)$. Se dice que $\{T_\alpha\}_{\alpha \in A}$ es \textbf{puntualmente acotada} si
    \begin{equation*}
        \forall x \in X : \{T_\alpha(x)\}_{\alpha \in A} \subseteq Y \text{ está acotada}
    \end{equation*}
\end{definicion}

\observacion{
    Si $\{T_\alpha\}$ está acotada uniformemente, entonces $\{T_\alpha\}$ está acotada puntualmente. De hecho:
    \begin{equation*}
        \|T_\alpha(x)\| \leq \|T_\alpha\| \cdot \|x\| \leq M \|x\| \implies \{T_\alpha(x)\} \subseteq Y \text{ está acotada}
    \end{equation*}
}

\begin{teo}[Teorema de Banach-Steinhaus o Principio de Acotación Uniforme]
    Sea $X$ un espacio de Banach e $Y$ un espacio normado, entonces toda familia $\{T_\alpha\}_{\alpha \in A} \subseteq BL(X, Y)$ puntualmente acotada es también uniformemente acotada. \\
    \noindent
    Es decir, si $\forall x \in X$ $\{T_\alpha(x)\}_{\alpha \in A} \subseteq Y$ está acotada, entonces 
    $$\exists M > 0 : \|T_\alpha\|_{BL(X, Y)} \leq M \quad \forall \alpha \in A$$

\begin{proof}
    Definimos $X_n := \{ x \in X \mid \|T_\alpha(x)\| \leq n \quad \forall \alpha \in A \} \quad \forall n \in \mathbb{N}$, para saber que es cerrado definimos para cada índice $\alpha \in A$ la aplicación
    \Func{h_\alpha}{X}{\R}{x}{\|T_\alpha(x)\|}
    continua por ser la composición de $T_\alpha$ con la norma de $Y$, que es continua. Tenemos entonces que 
    $$\{ x \in X \mid \|T_\alpha(x)\| \leq n \} = h_\alpha^{-1}([0, n])$$
    cerrado por ser la preimagen de un conjunto cerrado bajo una aplicación continua, y usando la intersección infinita tenemos
    $$X_n = \bigcap_{\alpha \in A} \{ x \in X \mid \|T_\alpha(x)\| \leq n \} = \bigcap_{\alpha \in A} X_{n, \alpha}$$
    cerrado por ser la intersección de cerrados. Además, se cumple 
    $$X = \bigcup_{n \in \mathbb{N}} X_n$$
    porque $\{T_\alpha\}_{\alpha \in A}$ es puntualmente acotada. Por el Teorema de Baire \ref{teo:baire}, $X$ es de II categoría ($X$ no es unión numerable de conjuntos raros), por lo tanto $\exists n_0 \in \mathbb{N} : X_{n_0}$ no es raro, es decir, $\exists B(x_0, r) \subseteq X_{n_0}$ con $x_0\in X$ y $r>0$, de lo cual
    \begin{align*}
        \forall z \in X : \|z\| \leq 1, \quad \forall \alpha \in A, \quad  \dfrac{r}{2} \|T_\alpha(z)\| &= \|T_\alpha\left(\dfrac{r}{2}z\right)\| \leq \|T_\alpha\left(\dfrac{r}{2}z + x_0\right)\| + \|T_\alpha(-x_0)\| \leq \\
        &\overset{(*)}{\leq} n_0 + |-1| \ \|T_\alpha(x_0)\| \overset{(**)}{\leq} 2n_0
    \end{align*}
    donde en $(*)$ hemos usado que $rz + x_0 \in B(x_0, r) \subseteq X_{n_0}$, es decir, $\|T_\alpha(rz + x_0)\| \leq n_0$ para todo $\alpha \in A$, y en $(**)$ hemos usado que $x_0\in X_{n_0}$. Por tanto
    \begin{equation*}
        \|T_\alpha(z)\| \leq \frac{2n_0}{r} \quad \forall z : \|z\| \leq 1 \quad \text{y por tanto} \quad \|T_\alpha\| = \sup_{\|z\|=1} \|T_\alpha(z)\| \leq \frac{n_0}{r} 
    \end{equation*}
\end{proof}
\end{teo}

\begin{notacion}
    Notar que a veces aludiremos a dicho teorema con la abreviación \tbf{P.A.U.} (Principio de Acotación Uniforme).
\end{notacion}

\begin{coro}
    Sea $X$ un espacio de Banach e $Y$ un espacio normado, si $\{T_n\}_{n \in \mathbb{N}} \subseteq BL(X, Y)$ converge puntualmente, es decir
    \begin{equation*}
        T_n(x) \longrightarrow T(x) \quad \forall x \in X \text{ en } Y
    \end{equation*}
    entonces $T$ es lineal y continuo.
\begin{proof}
    Dado $x \in X$, $(T_n(x))_n$ converge puntualmente $T(x)$, por lo tanto $(T_n(x))_n$ está acotada. Por el Teorema de Banach-Steinhaus, $\{T_n\}_n$ es uniformemente acotada, es decir
    $$\exists M > 0 : \|T_n\| \leq M \quad \forall n$$
    Por lo tanto, $\|T_n(x)\| \leq M \cdot \|x\|$ para todo $n$, y por lo tanto $\|T(x)\| \leq M \cdot \|x\|$. Esto implica que $T$ es acotado. Además, dado $x, y \in X$ y $\lambda \in \mathbb{K}$, tenemos que
    $$\|T(x + y) - T(x) - T(y)\| = \lim_{n} \|T_n(x + y) - T_n(x) - T_n(y)\|=0$$
    y
    $$\|T(\lambda x) - \lambda T(x)\| = \lim_{n} \|T_n(\lambda x) - \lambda T_n(x)\| = 0$$
    por lo tanto $T$ es lineal, y como hemos visto antes que era acotado, entonces también es continuo al ser lineal.
\end{proof}
\end{coro}

\begin{coro}
    Si $X$ es un espacio de Banach e $Y$ normado, entonces 
    $$BL(X, Y) \subseteq L(X, Y) = \{ T: X \to Y \text{ lineales} \}$$ 
    es cerrado.
\begin{proof}
    Sea $(T_n) \subseteq BL(X, Y)$ tal que $T_n \longrightarrow T$ en $L(X, Y)$, es decir, $T_n(x) \longrightarrow T(x)$ para todo $x \in X$. Por el Corolario anterior, $T$ es lineal y continuo, es decir, $T \in BL(X, Y)$. Por lo tanto, $BL(X, Y)$ es cerrado en $L(X, Y)$.
\end{proof}
\end{coro}


\subsection{Aplicaciones}

\begin{definicion}
    Una sucesión $(x_n)_{n\in \mathbb{N}}$ de un espacio normado $X$ es \textbf{infinitesimal} si $x_n \to 0$ en $X$.
\end{definicion}

\noindent Sea $(a_n) \in \mathbb{R}$ tal que $\sum\limits_{k=1}^{\infty} a_k x_k$ converge en $\mathbb{R} \quad \forall \{x_n\} \in C_0$ donde 
\begin{equation*}
    C_0 = \{ \text{sucesiones } (x_k) \in \mathbb{R} \text{ infinitesimales: } x_k \to 0 \text{ con norma medida en } \ell_\infty \}
\end{equation*}
entonces $(a_n) \in \ell_1$, veámoslo. Sea 
$$F_n(x) = \sum_{k=1}^{n} a_k x_k\quad \text{con } x = (x_k) \in C_0$$
entonces: $F_n: C_0 \to \mathbb{R}$ es un funcional lineal y continuo. Veámos ambas cosas
\begin{itemize}
    \item \tbf{Linealidad}: Dado $x, y \in C_0$ y $\lambda \in \mathbb{R}$, se tiene que
    \begin{align*}
        F_n(x + y) &= \sum_{k=1}^{n} a_k (x_k + y_k) = \sum_{k=1}^{n} a_k x_k + \sum_{k=1}^{n} a_k y_k = F_n(x) + F_n(y) \\
        F_n(\lambda x) &= \sum_{k=1}^{n} a_k (\lambda x_k) = \lambda \sum_{k=1}^{n} a_k x_k = \lambda F_n(x)
    \end{align*}
    \item \tbf{Continuidad}: Dado $x \in C_0$, se tiene que $x_k \to 0$ y por lo tanto $x_k$ es acotada, ya que por propia definición del límite tenemos que 
    \begin{equation*}
        \forall \varepsilon > 0, \exists N \in \mathbb{N} : |x_n-0| < \varepsilon \quad \forall N \geq n
    \end{equation*} 
    y tomando $\varepsilon = 1$ se tiene que $|x_n| < 1$ para todo $n \geq N_0$, y como los términos anteriores a $N_0$ son finitos, entonces $|x_k| \leq M$ para todo $k$. Tenemos por ser acotada que  
    $$\|x\|_{\infty} := \sup_{k} |x_k| < \infty$$ 
    Por lo tanto, 
    \begin{equation}\label{eq:Fn-continuo}
        |F_n(x)| = \left|\sum_{k=1}^{n} a_k x_k\right| \leq \sum_{k=1}^{n} |a_k| \cdot |x_k| \leq \sup_{k} |x_k| \cdot \sum_{k=1}^{n} |a_k| = \|x\|_{\infty} \cdot \sum_{k=1}^{n} |a_k| < \infty
    \end{equation}
    para todo $n$, lo que implica que $F_n$ es acotado, y por tanto continuo para todo $n$.
\end{itemize}
\noindent
Tenemos por tanto que 
\begin{equation*}
    F_n\in C_0^*=\{T:C_0\to \R : T \text{ lineal y continuo}\}
\end{equation*}
donde $C_0^*$ es el espacio dual continuo de $C_0$.
\vspace{0.3cm}

\noindent
Observemos que $\|F_n\| = \sum_{k=1}^{n} |a_k| \quad \forall n$, para ello tenemos que tener presente la definición de norma para un funcional, es decir
\begin{equation*}
    \|F_n\| = \sup_{\|x\|_{\infty} = 1} |F_n(x)| 
\end{equation*}
y usando la ecuación \ref{eq:Fn-continuo} se tiene que $|F_n(x)| \leq \|x\|_{\infty} \cdot \sum_{k=1}^{n} |a_k|$, por lo tanto $\|F_n\| \leq \sum_{k=1}^{n} |a_k|$. Buscamos entonces una sucesión $x$ con $\|x\|_{\infty} = 1$ tal que $|F_n(x)| = \sum_{k=1}^{n} |a_k|$. \\ \\
\noindent
Dado $n\in \N$, si tomamos $\tilde{x} = (\tilde{x}_k)$ tal que $\tilde{x} := (1, 1, \dots, \overbrace{1}^{n}, 0, \dots, 0, \dots)$ entonces $\|\tilde{x}\|_{\infty} = 1$, y tenemos
\begin{equation}\label{eq:igualdad-norma-Fn}
    |F_n(\tilde{x})| = \left|\sum_{k=1}^{n} a_k \cdot 1\right| = \sum_{k=1}^{n} |a_k| \implies \|F_n\| = \sum_{k=1}^{n} |a_k|
\end{equation}
\\
\noindent
Por otra parte, la hipótesis indicaba que $\sum\limits_{k=1}^{\infty} a_k x_k$ converge para todo $x \in C_0$, por lo que sus sumas parciales están obligatoriamente acotadas. Como $F_n(x)$ es precisamente la suma parcial hasta $n$, entonces para cada $x$ existe una constante $C_x$ tal que
$$|F_n(x)| \leq C_x \quad \forall n \in \mathbb{N}$$
lo que significa que $\{F_n\}_n$ es puntualmente acotada.  \\

\noindent % //TODO: demostrar que C_0 es espacio de banach
Al saber que $C_0$ es un espacio de Banach, podemos aplicar el Teorema de Banach-Steinhaus, y por lo tanto $\{F_n\}_n$ es uniformemente acotada, es decir
$$\exists M > 0 :\|F_n\| \leq M \ \forall n\in \N \implies \sum_{k=1}^{n} |a_k| \leq M \quad \forall n\in \N \implies \sum_{k=1}^{\infty} |a_k| \leq M$$

\begin{coro}
    Sea $X$ un espacio normado y sea $B \subseteq X$, tenemos que
    \begin{equation*}
        \text{si } f(B)=\{f(x) \mid x \in B\} \text{ acotado } \forall f \in X^* \implies B \text{ está acotado}
    \end{equation*}
\begin{proof}
    Sea $J: X \to X^{**}$ el mapa canónico, es decir
    
    $$J_x(f) = f(x) \quad \forall f \in X^*, \forall x \in X$$
    Entonces $J$ es una isometría (lo vimos cuando se definió el espacio dual), es decir
    $$\|J_x\|_{X^{**}} = \|x\|_X \quad \forall x \in X$$
    \noindent
    Para que $(J_x)_{x\in B}$ esté acotada puntualmente, necesitamos que 
    \begin{equation*}
        \forall f\in X^* : \{J_x(f)\}_{x\in B} = \{f(x)\}_{x\in B} \text{ esté acotada}
    \end{equation*}
    pero como sabemos que $f(B) = \{f(x) \mid x \in B\}$ está acotado para todo $f \in X^*$ (hipótesis), entonces se tiene. \\ \\
    \noindent
    Por el Teorema de Banach-Steinhaus, que lo podemos aplicar ya que $J_x:X^* \to \R$ y como $\R$ espacio banach, y $X$ espacio normado, entonces $X^*$ es espacio de banach, entonces tenemos que $(J_x)_{x \in B}$ está acotada uniformemente en $X^{**}$, es decir:
    \begin{equation*}
        \exists M > 0 : \|J_x\|_{X^{**}} = \|x\|_X \leq M \quad \forall x \in B
    \end{equation*}
    lo que implica que $B$ está acotado.
\end{proof}
\end{coro}

\begin{ejercicio}
    Sea $a_n \in \mathbb{R}$ tal que $\sum_{k} a_k x_k < +\infty \quad \forall x = (x_k) \in \ell_1$, demuestra que entonces $(a_n) \in \ell_\infty$. \\ \\
    \noindent
    Sea 
    \Func{F_n}{\ell_1}{\R}{x}{\sum\limits_{k=1}^{n} a_k x_k}
    donde recordamos que 
    $$\ell_1=\{x=(x_k)_{k\in\N}\subseteq \R : \sum\limits_{k=1}^\infty |x_k| < +\infty\}$$
    con la norma $\|x\|_1 = \sum\limits_{k=1}^{\infty} |x_k|$, entonces $F_n$ es lineal y continuo, veámoslo, fijando $n\in\N$.
    \begin{itemize}
        \item \textbf{Linealidad}: Dado $x, y \in \ell_1$ y $\lambda \in \mathbb{R}$, se tiene que
        \begin{align*}
            F_n(x + y) &= \sum_{k=1}^{n} a_k (x_k + y_k) = \sum_{k=1}^{n} a_k x_k + \sum_{k=1}^{n} a_k y_k = F_n(x) + F_n(y) \\
            F_n(\lambda x) &= \sum_{k=1}^{n} a_k (\lambda x_k) = \lambda \sum_{k=1}^{n} a_k x_k = \lambda F_n(x)
        \end{align*}
        \item \textbf{Continuidad}: Dado $x \in \ell_1$, se tiene que 
        \begin{equation*}
            \|x\|_1=\sum_{k=1}^{\infty} |x_k| < +\infty
        \end{equation*}
        entonces, el objetivo será ver si $|F_n(x)|\leq M\|x\|_1$ con $M>0$, veámoslo
        \begin{equation*}
            |F_n(x)| = \left|\sum_{k=1}^{n} a_k x_k\right| \leq \sum_{k=1}^{n} |a_k| \cdot |x_k| \leq \max_{k} |a_k| \cdot \sum_{k=1}^{n} |x_k| \leq \max_{k} |a_k| \cdot \|x\|_1 < +\infty
        \end{equation*}
        por lo que $F_n$ es acotado, y por tanto continuo para todo $n$.
    \end{itemize}
    Observemos ahora que $\|F_n\|=\max\limits_{k=1,\dots, n}|a_k|$, para ello usaremos la definición de norma de operadores lineales, es decir
    \begin{equation*}
        \|F_n\|=\sup_{\|x\|_1}|F_n(x)| \leq \sup_{\|x\|_1} \max_{k=1,\dots,n} |a_k| \cdot \|x\|_1 = \max_{k=1,\dots,n} |a_k|
    \end{equation*}
    Viendo la desigualdad dada, buscaremos un $x\in \ell_1$ con $\|x\|_1=1$ tal que $|F_n(x)| = \max_{k=1,\dots,n} |a_k|$, y tendríamos la igualdad buscada. Sea $\tilde{x}\in \ell_1$ tal que 
    \begin{equation*}
        \tilde{x} = 
        \begin{cases} 
            0 & k \neq \tilde{k} \\ 
            a_{\tilde{k}} & k = \tilde{k} 
        \end{cases}
        \quad \text{ con } a_{\tilde{k}} = \max_{k=1,\dots,n} |a_k|
    \end{equation*}
    entonces tenemos que 
    $$|F_n(\tilde{x})| = \left|\sum_{k=1}^{n} a_k \tilde{x}_k\right| = |a_{\tilde{k}}| = \max_{k=1,\dots,n} |a_k|$$
    y por lo tanto $\|F_n\| = \max\limits_{k=1,\dots,n} |a_k|$. \\ \\
    \noindent
    Por otra parte, como bien demostramos que para todo $x\in \ell_1$ se tiene que $|F_n(x)|$ está acotada, entonces tenemos que $\{F_n\}_n$ es puntualmente acotada. Por el Teorema de Banach-Steinhaus, tenemos $\ell_1$ espacio de Banach, $\{F_n\}_n$ es uniformemente acotada, es decir
    \begin{equation*}
        \exists M > 0 : \|F_n\| = \max_{k=1,\dots,n} |a_k| \leq M \quad \forall n \in \mathbb{N} \implies |a_k| \leq M \quad \forall k \in \mathbb{N} \implies (a_n) \in \ell_\infty
    \end{equation*}
\end{ejercicio}

\begin{ejercicio}
    Sea $X = \{ \text{polinomios}, t \in [0, 1] \}$ con $\|p\| = \max\limits_{j=0, \dots, \text{deg } p} |a_j|$. Sea 
    \Func{F_n}{X}{\R}{\sum\limits_{j=1}^{n} a_j t^j}{\sum\limits_{j=1}^{n} a_j}
    mostrar que $F_n$ son continuos pero no uniformemente acotados. \\ \\
    \noindent
    Para demostrar que $F_n$ lo que queremos es ver que es un operador lineal y después demostrar la acotación, que equivale en operadores linales a la continuidad. Para la linealidad, dado $p, q \in X$ y $\lambda \in \mathbb{R}$, se tiene que
    \begin{align*}
        F_n(p + q) &= \sum_{j=1}^{n} (a_j + b_j) = \sum_{j=1}^{n} a_j + \sum_{j=1}^{n} b_j = F_n(p) + F_n(q) \\
        F_n(\lambda p) &= \sum_{j=1}^{n} \lambda a_j = \lambda \sum_{j=1}^{n} a_j = \lambda F_n(p)
    \end{align*}
    para la continuidad, dado $p \in X$, se tiene que
    \begin{equation*}
        |F_n(p)| = \left|\sum_{j=1}^{n} a_j\right| \leq \sum_{j=1}^{n} |a_j| \leq n \cdot \max_{j=1,\dots,n} |a_j| \overset{(*)}{\leq} n \cdot \|p\|
    \end{equation*}
    donde en $(*)$ hemos usado que la definición de norma toma el máximo de $|a_j|$, pero incluyendo $j=0$, de ahí que pongamos el $\leq$ en lugar de la igualdad. Al igual que antes aplicando la definición de norma de operadores lineales, se tiene que
    \begin{equation*}
        \|F_n\| = \sup_{\|p\|=1} |F_n(p)| \leq \sup_{\|p\|=1} n \cdot \|p\| = n
    \end{equation*}
    por lo que $F_n$ es continuo para todo $n$. \\ \\
    \noindent
    Siguiendo la misma dinámica de antes, definimos un polinomio concreto para ver la igualdad en $\|F_n\|$, es decir, definimos $p^*=\sum\limits_{k=1}^{n} t^k$, entonces
    \begin{equation*}
        |F_n(p^*)| = \left|\sum_{k=1}^{n} 1\right| = n \implies \|F_n\| = n
    \end{equation*}
    es claro que tenemos que $\{F_n\}$ no está uniformemente acotada, ya que $\|F_n\| = n$ para todo $n$, y por lo tanto no existe $M > 0$ tal que $\|F_n\| \leq M$ para todo $n$. De hecho, $(X, \|\cdot\|)$ no es un espacio de Banach y, por lo tanto, el P.A.U. no es válido. \\ \\
    \noindent
    Es importante notar aqui que la elección de $p^*$ esta hecha en función de $F_n$, es decir, el grado de $p^*$ es $n_0$ concreto para $F_{n_0}$, y por lo tanto para la \underline{acotación uniforme} no es un polinomio fijo para todos los $F_n$, sino que cada $F_n$ tiene su propio polinomio $p$ con el que se alcanza la igualdad en la norma de operadores lineales. Esto es de vital importancia, puesto que \tbf{SI} cumple la \underline{acotación puntual}, ya que si tomo mi $p^*$ particular para un $F_{n_0}$ concreto, entonces $|F_{n_0}|=n_0$, pero para $n>n_0$ tendríamos el mismo valor $|F_n(p^*)|=n_0$, por lo que la acotación puntual se cumple, al tener una sucesión convergente en la constante $n_0$.
\end{ejercicio}
